\section{Definir: punto de vista y principios}
\label{sec:definir}

\subsection{Punto de vista}

\begin{nota}
Martina, que quiere usar menos el celular pero abandona las apps que lo intentan, \textbf{necesita} que desconectarse le dé algo concreto a cambio sin tener que acordarse de nada, \textbf{porque} las herramientas actuales le piden esfuerzo y solo le devuelven estadísticas.
\end{nota}

\subsection{¿Cómo podríamos...?}

\begin{enumerate}[label=\textbf{HMW\arabic*}]
\item ¿Cómo podríamos registrar una pausa sin que el usuario haga nada?
\item ¿Cómo podríamos lograr que una hora sin el celular valga algo fuera de la pantalla?
\item ¿Cómo podríamos evitar que alguien sume puntos durmiendo, sin que la regla se sienta injusta?
\item ¿Cómo podríamos canjear un premio en el mostrador en segundos, incluso sin señal?
\end{enumerate}

\subsection{Principios de diseño}

De esas preguntas salen cinco principios. Los usamos para decidir cada pantalla y cada requisito.

\begin{table}[H]
\centering\small
\begin{tabularx}{\textwidth}{@{}l l L@{}}
\toprule
& \textbf{Principio} & \textbf{Qué implica} \\
\midrule
P1 & Cero toques para lo central & La pausa se detecta sola. El usuario no inicia ni termina nada. \\
P2 & La moneda es el tiempo & 1 minuto de pausa es 1 punto. Los precios también se muestran en horas (“te falta 1 h”). \\
P3 & Una acción, al alcance del pulgar & Una sola acción principal por pantalla, grande y abajo (ley de Fitts). \\
P4 & Reglas a la vista & Si una pausa no suma, la app dice por qué. Nada de reglas ocultas. \\
P5 & Sin red también anda & Detectar, sumar y consultar funcionan offline. El canje queda pendiente, no falla. \\
\bottomrule
\end{tabularx}
\end{table}
